\documentclass[a4paper, 12pt]{article}

\usepackage[utf8]{inputenc}
\usepackage[T1]{fontenc}
\usepackage{textcomp}
\usepackage{amssymb}
\usepackage{newtxtext} \usepackage{newtxmath}
\usepackage{amsmath, amssymb}
\newtheorem{problem}{Problem}
\newtheorem{example}{Example}
\newtheorem{lemma}{Lemma}
\newtheorem{theorem}{Theorem}
\newtheorem{problem}{Problem}
\newtheorem{example}{Example} \newtheorem{definition}{Definition}
\newtheorem{lemma}{Lemma}
\newtheorem{theorem}{Theorem}
\DeclareMathAlphabet{\mathcal}{OMS}{cmsy}{m}{n}

\begin{document}

En sólo dos horas volveré a dormir y, si mi preparación sirvió de algo, no
volveré a despertar. No estaré muerto, desde luego, y con cierta laxitud podría
incluso decirse que viviré toda una eternidad, o lo más cercano a la eternidad
que una consciencia puede barruntar.

~

Hace dos años ya que terminé mi formación en la Universidad de Pennsylvania.
Allí estudié cuatro años el sueño de ondas lentas, la sincronización
espacio-temporal de ciertas frecuencias del sueño, y, lo que es más importante,
la desactivación de las regiones del cerebro que sirven de interfaz para la
metacognición. Despiertos decimos: «estoy pensando»; dormidos, casi nunca
podemos decir «estoy soñando», y tenemos un conocimiento relativamente claro de
qué poblaciones de neuronas, activas en la vigilia, se inhiben en el sueño,
dándonos y quitándonos alternativamente el don de la auto-referencialidad. 

~ 

Cuando abandoné el laboratorio, Paulina ya había muerto y yo estaba destinado a
doctorarme en Ginebra. Volví a Buenos Aires sólo como un acto silencioso de
protesta: de protesta contra la vida, que me había dado tanto mientras
abandonaba a Paulina, y de protesta contra la muerte, que yo empezaba a anhelar
como un alivio. Dejé mi carrera
























\end{document}



